\chapter{Introduction}
\label{chap:intro}

\section{Background and Motivation}

Abdominal CT scans are among the most frequently ordered imaging studies in clinical practice.
Radiologists rely on them for cancer staging, surgical planning, post-operative follow-up, and
emergency diagnosis. At the heart of many of these decisions lies a deceptively simple question:
which voxels belong to which organ? Answering it accurately and quickly matters. An error in
liver boundary delineation before hepatic resection, for instance, can translate directly into
operative risk.

Manual organ delineation is tedious. For a single CT volume containing $200$--$400$ axial slices
and 13 structures, an experienced radiologist may spend anywhere from one to three hours to produce
a high-quality annotation. At scale --- across tens of thousands of scans in a busy hospital ---
that cost is prohibitive. Automated segmentation reduces turnaround time from hours to seconds
while removing the inter-observer variability that plagues manual methods.

Deep convolutional architectures have reshaped the segmentation landscape. The encoder--decoder
design popularised by U-Net~\cite{ronneberger2015} elegantly captures both local texture and global
context through skip connections, and subsequent work has extended it with attention mechanisms,
volumetric processing, and self-supervised pretraining. Yet gaps remain: 3D methods are
memory-hungry, transformer-based models require enormous training sets, and many published results
still rely on per-slice evaluation metrics that overstate true performance. This project addresses
those gaps directly through a carefully implemented 2D Attention U-Net evaluated under a rigorous
per-volume protocol.

\section{Problem Statement}

The task is \emph{multi-organ semantic segmentation}: given an axial CT slice, assign each pixel
one of 14 labels (background plus 13 abdominal organs). The challenge compounds three
difficulties: extreme class imbalance (the pancreas and adrenal glands collectively occupy less
than 1\% of foreground voxels), high morphological variability across patients, and the 2D
processing constraint imposed by GPU memory limits on freely available hardware.

A secondary but equally important problem is evaluation integrity. Early experiments revealed that
computing Dice per slice and averaging inflates scores significantly --- empty slices where both
prediction and ground-truth are zero contribute a spurious Dice of 1.0 under na\"\i{}ve
implementations. The corrected per-volume protocol used throughout this report closes that
loophole.

\section{Objectives}

\begin{itemize}
    \item Implement a 2D Attention U-Net with attention gates on all skip connections for
          multi-organ abdominal CT segmentation.
    \item Train and evaluate on the FLARE22 benchmark using a per-volume evaluation protocol that
          matches the challenge's own scoring conventions.
    \item Quantify the contribution of attention gates, Test-Time Augmentation, and the cosine
          learning-rate schedule through a controlled ablation study.
    \item Demonstrate competitive Dice and HD95 results relative to published nnU-Net 3D baselines,
          using only 2D processing and a single T4 GPU.
\end{itemize}

\section{Contributions}

The principal contributions of this project are summarised in Table~\ref{tab:contributions}.

\begin{table}[h!]
\centering
\caption{Summary of contributions.}
\label{tab:contributions}
\begin{tabular}{p{3.5cm} p{7cm} p{2.5cm}}
\toprule
\textbf{Contribution} & \textbf{Description} & \textbf{Chapter} \\
\midrule
Attention U-Net implementation &
  Full PyTorch implementation with attention gates on all four skip connections, Dropout2d
  bottleneck, and AMP-enabled training. &
  Chapter~\ref{chap:methodology} \\[6pt]
Per-volume evaluation fix &
  Identified and corrected the per-slice Dice inflation bug; established a per-volume evaluation
  loop consistent with FLARE22 conventions. &
  Chapter~\ref{chap:eval} \\[6pt]
TTA ensemble &
  Four-way horizontal/vertical flip ensemble with softmax averaging, adding $+$0.06 Dice at zero
  extra training cost. &
  Chapter~\ref{chap:eval} \\[6pt]
Ablation study &
  Controlled comparison across Runs~014--017 isolating the effect of data volume, attention gates,
  cosine LR, and TTA. &
  Chapter~\ref{chap:ablation} \\[6pt]
Competitive results on tight hardware &
  Mean Dice 0.8769, HD95 7.59~mm on 13 organs; numerically exceeding the nnU-Net 3D baseline
  on 7 of 13 organs using only 2D slices and a single T4 GPU. &
  Chapter~\ref{chap:results} \\
\bottomrule
\end{tabular}
\end{table}

\section{Report Organisation}

Chapter~\ref{chap:lit} reviews the literature on medical image segmentation, attention mechanisms,
and the FLARE22 benchmark. Chapter~\ref{chap:data} describes the dataset and preprocessing
pipeline. Chapter~\ref{chap:methodology} details the Attention U-Net architecture and loss
function. Chapter~\ref{chap:train} covers the training configuration. Chapter~\ref{chap:eval}
defines the evaluation protocol. Chapter~\ref{chap:results} presents the quantitative results.
Chapter~\ref{chap:ablation} reports the ablation study and discussion. Chapters~\ref{chap:conc}
and \ref{chap:future} give the conclusion and directions for future work.
